% === Begin Label Data ===
% ID: 87
% 难度星级: 
% 标签: 
% 备注: 
% 组卷引用次数: 0
% === End  Label Data ===

\begin{problem}{2024}{G}{新课标I卷}{19}{数列}
设 $m$ 为正整数，数列 $a_1, a_2, \cdots, a_{4m+2}$ 是公差不为 0 的等差数列，若从中删去两项 $a_i$ 和 $a_j(i<j)$ 后，剩余的 $4m$ 项可以平均分成 $m$ 组，且每组的 4 个数都能构成等差数列，则称数列 $a_1, a_2, \cdots, a_{4m+2}$ 是 $(i,j)-$ 可分数列.

(1) 写出所有的 $(i,j)$，$1\leqslant i<j \leqslant 6$，使得 $a_1, a_2, \cdots, a_6$ 是 $(i,j)-$ 可分数列;

(2) 当 $m\geqslant 3$ 时，证明：数列 $a_1, a_2, \cdots, a_{4m+2}$ 是 $(2,13)-$ 可分数列;

(3) 从 $1, 2, \cdots, 4m+2$ 中一次任取两个数 $i$ 和 $j(i<j)$，记数列 $a_1, a_2, \cdots, a_{4m+2}$ 是 $(i,j)-$ 可分数列的概率为 $P_m$，证明：$P_m>\dfrac{1}{8}$.
\end{problem}